% Generated shared reviewed material.
% Full derivation of the scoped arithmetic-family moment and masked remainder.
\section{Family identification, the approximate functional equation and masks}
\label{mm:details}
We prove the analytic statements of Section~\ref{mm:core} from the
explicit classical inputs stated there. Write $I(S)$ for the nonzero
integral ideals prime to $S$, and $\mathcal O_S$ for the ring of
$S$-integers. Constants may depend on the fixed field, $S$, $\nu$ and
any declared positive small exponent. No such constant is asserted
uniform in the conductor or infinity type of a varying $\nu$.

\subsection{A finite comparison with the literal Kummer rows}
Use Appendix~II of BGL, with $n=6$. Its fixed local modulus and
representatives may be chosen once, since $S$ contains the primes above
$6$ and $\mathcal O_S$ has class number one. For each
$\mathfrak a\in I(S)$ the construction writes
\[
 \mathfrak a=(x)E\mathfrak g^6,\qquad
 \alpha_{\mathfrak a}=xm_E,\qquad
 (m_E)_S=E,
\]
so $(\alpha_{\mathfrak a})_S=\mathfrak a\mathfrak g^{-6}$.
Choose $y\in K^\times$ generating $\mathfrak g$ as an
$\mathcal O_S$-ideal. If $(u)_S=\mathfrak a$, then
\[
 e=\alpha_{\mathfrak a}y^6/u\in\mathcal O_S^\times,\qquad
 [\alpha_{\mathfrak a}]=[eu]\quad\hbox{in }K^\times/K^{\times6}.
\]
There is a precision here: Appendix~II allows choosing $\mathfrak g$
relative to a tested denominator. Fix one admissible global decomposition
first. Its radicand gives the canonical character
$\vartheta_{\mathfrak a}$ on all primes outside the finite set
$S\cup\operatorname{supp}(\mathfrak a\mathfrak g)$ where that
choice applies. Uniqueness of Hecke characters, equivalently Chebotarev
in the finite abelian quotient, then identifies the underlying primitive
characters everywhere. With matching arithmetic-Frobenius conventions,
\begin{equation}\label{mm:unit-comparison}
 \kappa_u=\kappa_{e^{-1}}\vartheta_{\mathfrak a}.
\end{equation}
If the residue-symbol convention is conjugated, conjugate the canonical
family and its coefficient sequence; the same large sieve follows.

The group $\mathcal O_S^\times/(\mathcal O_S^\times)^6$ is finite.
Thus~\eqref{mm:unit-comparison} uses one of a finite set of fixed
Kummer twists $\eta$, all unramified outside $S$ and of bounded
$S$-conductor. This is a comparison of primitive characters, not a
pointwise multiplication rule for independently zero-extended arrays at
bad primes. In particular, a conductor drop must not be replaced by a
product of artificial zeros.

Since $\mathcal O$ is principal with six units, for fixed
$\mathfrak a\in I(S)$ and an $S$-supported integral ideal
$\mathfrak b$ there are at most six $u$ satisfying
$(u)=\mathfrak a\mathfrak b$. For any nonnegative row quantity, the
literal moment is consequently at most
\begin{equation}\label{mm:literal-sum}
 6\sum_{\mathfrak b\ S\text{-supported}\atop N\mathfrak b\le H}
   \sum_{\eta}\sum_{\mathfrak a\in I(S)\atop N\mathfrak a\le H/N\mathfrak b}
 \left|L\bigl(s,\operatorname{prim}(\nu\eta
                       \vartheta_{\mathfrak a})\bigr)\right|^2.
\end{equation}
An ideal bound $X^{1+\varepsilon}T^{1+\varepsilon}$ therefore transfers
to all literal rows, since
\[
 \sum_{\mathfrak b\ S\text{-supported}}(N\mathfrak b)^{-1-\varepsilon}
 =\prod_{p\in S}(1-(Np)^{-1-\varepsilon})^{-1}<\infty.
\]
No quotient of the original unit rows or deletion of $S$-prime powers
has occurred.

\subsection{Fixed twists, exact tame conductors and the squarefree sieve}
Fix $\xi=\nu\eta$ and write
$\rho_{\mathfrak a}=\operatorname{prim}
(\xi\vartheta_{\mathfrak a})$. Its fixed conductor is supported in
$S$ before the Kummer twist. On columns prime to $S$, $|\xi(\mathfrak b)|=1$.
BGL Theorem~1.3 thus gives, for any fixed $\mathfrak a_1$,
\begin{align}\label{mm:ls}
 &\sum_{\mathfrak a_0\in I(S)\ \mathrm{squarefree}\atop N\mathfrak a_0\le M}
 \left|\sum_{\mathfrak b_1\in I(S)\ \mathrm{squarefree}\atop N\mathfrak b_1\le B}
 \lambda_{\mathfrak b_1}\xi(\mathfrak b_1)
 \vartheta_{\mathfrak a_1}(\mathfrak b_1)
 \vartheta_{\mathfrak a_0}(\mathfrak b_1)\right|^2\notag\\
 &\hspace{1cm}\ll (MB)^\varepsilon
 \{M+B+(MB)^{2/3}\}\sum_{\mathfrak b_1}|\lambda_{\mathfrak b_1}|^2.
\end{align}
The first fixed multiplier preserves the coefficient norm; the second
has modulus at most one. Restricting the positive outer row sum is
harmless. A coprimality condition with another fixed ideal belongs to
$\lambda$. In particular, $\nu$ need not have order dividing six and
is not asserted to belong to BGL's finite ray-class character set.

There is a unique valuation decomposition
\[
 \mathfrak a=\mathfrak a_0\mathfrak a_1\mathfrak a_2^6,
 \quad \mathfrak a_0\ \hbox{squarefree},\quad
 \mathfrak a_1\ \hbox{squarefull and sixth-power-free},\quad
 (\mathfrak a_0,\mathfrak a_1)=1.
\]
Indeed, for $v_p(\mathfrak a)=6q+r$, $0\le r\le5$, put the residual
prime in $\mathfrak a_0$ if $r=1$, and its exponent $r$ in
$\mathfrak a_1$ if $2\le r\le5$. BGL Lemma~2.1 gives
$\vartheta_{\mathfrak a\mathfrak b^6}=\vartheta_{\mathfrak a}$ and
$\vartheta_{\mathfrak a_0\mathfrak a_1}
=\vartheta_{\mathfrak a_0}\vartheta_{\mathfrak a_1}$ on ideals
prime to $S$ in this sixth-power-free case. Those are the identities
needed here; no arbitrary multiplicativity assertion is substituted.

At a prime outside $S$, tame Kummer inertia at residual exponent $r$
has order $6/\gcd(6,r)$. It is nontrivial for $r=1,\ldots,5$ and
has conductor exponent one. The fixed $\xi$ is unramified there, so
cannot cancel that inertia. At primes in $S$ the local sixth-power
classes are finite, hence only finitely many primitive conductor types
occur after multiplication by the fixed $\xi$. Consequently
\begin{equation}\label{mm:conductor}
 Q_{\rho_{\mathfrak a}}=|D_K|N\mathfrak f_{\rho_{\mathfrak a}}
 \le C_{\xi,S}N(\mathfrak a_0\mathfrak a_1),
\end{equation}
independently of $\mathfrak a_2$. The equality of the conductor's
outside-$S$ part with $\operatorname{rad}(\mathfrak a_0\mathfrak a_1)$
does not identify that conductor with the original larger deleted set.

\subsection{Removing the row-dependent functional-equation weights}
First take finite-order $\nu$. All characters then have trivial infinity
type at the complex place. The primitive completion may be normalized as
\[
 \Lambda(s,\rho)=Q_\rho^{s/2}\gamma_K(s)L(s,\rho),\qquad
 \gamma_K(s)=2(2\pi)^{-s}\Gamma(s),\qquad
 \Lambda(s,\rho)=\epsilon_\rho\Lambda(1-s,\bar\rho),
 \quad |\epsilon_\rho|=1.
\]
For nonprincipal rows the completed function is entire. Apply this
functional equation to the Mellin integral with even auxiliary factor
$G(z)=e^{z^2}$ and denominator $z$, shifting across $z=0$. At
$s=1/2+it$ this gives two weighted Dirichlet series, truncated below to
polynomials; the dual prefactor has modulus one. Principal rows are excluded from
this pole-free step and treated separately below.

For fixed $\mathfrak a_2$, let $X=H/(N\mathfrak a_2)^6\ge1$ and
$T=1+|t|$. An explicit normalized kernel is
\begin{equation}\label{mm:afe-kernel}
 V_t(x)=\frac1{2\pi i}\int_{(\eta)}
 e^{z^2}\frac{\Gamma(s+z)}{\Gamma(s)}T^{-z}x^{-z}\,\frac{dz}{z},
 \qquad x=\frac{2\pi N\mathfrak b}{\sqrt{Q_\rho}T},
\end{equation}
with $\eta>0$ small. Its factors combine to
$(\sqrt{Q_\rho}/(2\pi N\mathfrak b))^z$ in the actual completed-L
kernel; there is no extra factor of $T$. The denominator $z$ supplies
the residue at zero.
On $z=\eta+iy$, Stirling and $|t+y|\le T(1+|y|)$ give
\[
 \left|\frac{\Gamma(1/2+\eta+i(t+y))}
                 {\Gamma(1/2+it)T^{\eta+iy}}\right|
 \ll_\eta (1+|y|)^{C_\eta}e^{\pi|y|/2}.
\]
The factor $e^{-y^2}$ makes this an integrable majorant uniformly in
$t$. The analogous bound holds on every fixed positive shifted line
and for the dual gamma kernel. Moving far enough to the right gives
arbitrary power decay beyond a common total-ideal cutoff
\begin{equation}\label{mm:common-cutoff}
 B'=X^{1/2}T(HT)^\eta,
\end{equation}
up to fixed $\xi,S,K$ constants. For any requested $A>0$, a sufficiently
large positive kernel shift, depending also on $\eta$, makes the
truncation error $O((HT)^{-A})$ uniformly in all rows in this block.
Ideal counting $O_K(Y)$ bounds the absolute tails. Such errors are
negligible after the finite row sums by taking $A$ larger if needed.

Insert~\eqref{mm:afe-kernel} on $\operatorname{Re}z=\eta$.
The factor $(\sqrt{Q_\rho}T)^z$ is independent of the column and is a
row multiplier of modulus at most $(CXT^2)^{\eta/2}$.
Its imaginary part is a phase. The remaining column coefficient is
$(N\mathfrak b)^{-1/2-\eta-it-iy}$, apart from fixed constants.
Minkowski or Cauchy with the integrable gamma majorant therefore removes
the row-dependent weight at cost $(HT)^{O(\eta)}$ in the squared norm.
No independence of the root numbers is assumed or needed.

For each column use the unique decomposition
\[
 \mathfrak b=\mathfrak b_S\mathfrak b_1\mathfrak b_2,
 \quad \mathfrak b_S\ S\hbox{-supported},\quad
 \mathfrak b_1\ \hbox{squarefree},\quad
 \mathfrak b_2\ \hbox{squarefull},\quad
 (\mathfrak b_1,\mathfrak b_2)=1,
\]
where $\mathfrak b_1,\mathfrak b_2\in I(S)$; a prime of outside-$S$
valuation one goes in $\mathfrak b_1$, and a prime of valuation at least
two goes with its full exponent in $\mathfrak b_2$. For fixed
$\mathfrak b_S,\mathfrak b_2$ retain
$\rho_{\mathfrak a}(\mathfrak b_S)$ as the original primitive row
value. On the good factors the exact canonical identity applies.
The remaining factors involving $\mathfrak b_2$, including
$\vartheta_{\mathfrak a_0}(\mathfrak b_2)$, are bounded row
multipliers. Their moduli are at most one; only then are they discarded
in an inequality. In particular the value at $\mathfrak b_S$ is not
replaced by the product of zero-extended primitive arrays.

For the $\mathfrak b_1$ polynomial the coefficient sequence is
independent of $\mathfrak a_0$; it includes the indicator
$(\mathfrak b_1,\mathfrak b_2)=1$ and the common cutoff
$N\mathfrak b_1\le B'/N(\mathfrak b_S\mathfrak b_2)$.
Padding it by zeros up to $B'$ is legitimate. Its squared norm is at most
$\zeta_K(1+2\eta)$. Minkowski over $\mathfrak b_S,\mathfrak b_2$
is legitimate because
\[
 \sum_{\mathfrak b_S,\mathfrak b_2}
       N(\mathfrak b_S\mathfrak b_2)^{-1/2-\eta}<\infty.
\]
The finite Euler product handles $\mathfrak b_S$; the squarefull series
converges above $1/2$. Apply~\eqref{mm:ls} to each such polynomial and
to the conjugate dual family. For fixed $\mathfrak a_1,\mathfrak a_2$
the resulting energy is
\begin{equation}\label{mm:block-energy}
 (HT)^{O(\eta)}\{M+B'+(MB')^{2/3}\},\qquad
 M=X/N\mathfrak a_1.
\end{equation}
The additional condition $(\mathfrak a_0,\mathfrak a_1)=1$ simply
restricts the positive row sum. The large-sieve small exponent can be
chosen sufficiently small along with $\eta$ and absorbed in the
indicated small loss.

\subsection{All-ideal summation and principal rows}
Every squarefull ideal $\mathfrak c$ has a unique expression
$\mathfrak d^2\mathfrak e^3$ with $\mathfrak e$ squarefree.
This follows primewise by using exponent zero or one in $\mathfrak e$
according to the parity of an exponent at least two. Therefore
\[
 \#\{\mathfrak c\ \mathrm{squarefull}:N\mathfrak c\le X\}
 \ll_K X^{1/2},\qquad
 \sum_{\mathfrak c\ \mathrm{squarefull}}(N\mathfrak c)^{-s}<\infty
 \quad(s>1/2).
\]
For the count, sum
$O((X/(N\mathfrak e)^3)^{1/2})$ and use
$\zeta_K(3/2)<\infty$. Summing~\eqref{mm:block-energy} over
$\mathfrak a_1$ gives
\[
 (HT)^{O(\eta)}
 \{X+B'\sqrt X+(XB')^{2/3}\}
 \ll (HT)^{O(\eta)}XT,
\]
since the squarefull series converges at $1$ and $2/3$,
$B'=\sqrt X\,T(HT)^\eta$, and $T^{2/3}\le T$ for $T\ge1$.
The $\mathfrak a_2$ sum is bounded by
$H\sum_{\mathfrak a_2}(N\mathfrak a_2)^{-6}=H\zeta_K(6)$.
Choosing all internal losses small for the requested final
$\varepsilon$ proves the fixed-twist all-ideal central moment.

There are $O_K(H)$ ideal rows in total. On any row where the primitive
character is principal its $L$-function is $\zeta_K$, and ordinary
scalar convexity gives
$|\zeta_K(1/2+it)|^2\ll_\varepsilon T^{1+\varepsilon}$.
This treats those rows without a pole-free approximate functional
equation. Finally~\eqref{mm:literal-sum} proves~\eqref{mm:central}
for every actual nonzero Eisenstein element.

\subsection{Poisson transfer with an explicit constant}
For a nonprincipal row define
\[
 f_\psi(s)=\frac{L(s,\psi)}{(s+1)^{5/4}},
\]
using the holomorphic logarithm on $\operatorname{Re}(s+1)>0$.
On a principal row additionally multiply by $(s-1)/(s+1)$.
Analytic continuation and ordinary degree-two convexity, used separately
at each fixed conductor, show that $f_\psi$ is bounded analytic on the
right half-plane $\operatorname{Re}s>1/2$, with continuous boundary
and decay at infinity. The principal pole is removed before applying
Poisson representation to its real and imaginary parts. Jensen gives
\[
 |f_\psi(1/2+\delta+it)|^2
 \le\int_{\mathbb R}P_\delta(t-v)
             |f_\psi(1/2+iv)|^2\,dv.
\]
The exact row family is finite for fixed $H$, so its sum interchanges
with the integral without a limiting-family issue. The boundary
pole-removal factor has modulus at most one.

Take $0<\varepsilon_0\le1/4$ and
$\alpha=3/2-\varepsilon_0\in[5/4,3/2)$. Since
$|3/2+iv|\ge(1+|v|)/\sqrt2$, the central squared moment of the
$f_\psi$ is bounded by
$3C_0H^{1+\varepsilon_0}w_\alpha(v)$, where
$w_\alpha(v)=(1+|v|)^{-\alpha}$.
For $|t|\ge2$, split $|v-t|\le|t|/2$ from its complement. The
near contribution to $P_\delta*w_\alpha$ is at most
$2^\alpha w_\alpha(t)<3w_\alpha(t)$. In the far region use
$P_\delta(t-v)\le4\delta/(\pi t^2)$ and
$\int w_\alpha=2/(\alpha-1)$, giving at most
\[
 \frac{2\delta3^\alpha}{\pi(\alpha-1)}w_\alpha(t)
 \le8w_\alpha(t).
\]
Here $(1+|t|)^\alpha/t^2\le3^\alpha/4$ for $|t|\ge2$.
For $|t|\le2$, the convolution is at most one and hence at most
$3^\alpha w_\alpha(t)<6w_\alpha(t)$. Thus
$P_\delta*w_\alpha\le11w_\alpha$ uniformly for $0<\delta<1/2$.
Restoring the denominator costs less than
$6(1+|t|)^{5/2}$, and the inverse principal factor is bounded by
$(3/2+\delta)/(1/2-\delta)$. The product $11\cdot3\cdot6=198<200$
proves~\eqref{mm:poisson-constant}. The constant is uniform as
$\delta\downarrow0$, but not as the contour approaches the principal
pole line $\operatorname{Re}s=1$. No zero-free hypothesis appears.

\subsection{The separately scoped unitary infinity type}
Let $\nu(\mathfrak a)=\nu_0(\mathfrak a)(N\mathfrak a)^{-i\tau}$
be fixed and unitary, with angular parameter $m$ and gamma factor
$\gamma_\nu$ stated in Section~\ref{mm:core}. The external primitive
functional equation is explicitly
\[
 Q_\rho^{s/2}\gamma_\nu(s)L(s,\rho)
 =\epsilon_\rho Q_\rho^{(1-s)/2}
                 \gamma_{\bar\nu}(1-s)L(1-s,\bar\rho),
 \qquad |\epsilon_\rho|=1.
\]
A constant norm-twist phase in another completion convention is absorbed
in $\epsilon_\rho$. All finite Kummer twists share this infinity type.
Replace $T$ in the kernel by $T_\nu=1+|t+\tau|+|m|$.
The numerator gamma ratio now has real base $1/2+|m|/2>0$ and
vertical parameter $t+\tau$; its conjugate dual has parameter
$-(t+\tau)$. Stirling gives the same integrable Gaussian-majorized
kernel on each fixed shifted line. Constants can depend on $m,\tau$;
$T_\nu$ is comparable to $T$ for this one fixed $\nu$.
The finite-prime conductor analysis, modulus-one coefficient twist,
all-ideal decomposition and all original primitive row values remain
unchanged.

A pure norm row has $L(s,\rho)=\zeta_K(s+i\tau)$; exclude it from
the pole-free approximate functional equation and use scalar convexity
at $t+\tau$. In the Poisson step cancel its actual pole by
\[
 h_\tau(s)=\frac{s-1+i\tau}{s+1+i\tau}.
\]
This factor is analytic in the relevant right half-plane. On its
boundary $|h_\tau|\le1$, while inside
$|h_\tau^{-1}|\le(3/2+\delta)/(1/2-\delta)$.
Hence the same transfer factor $200$ applies to the now
$\nu$-dependent central constant. If the opposite norm convention is
chosen, replace $\tau$ by $-\tau$ in the gamma factor and pole alike.
This proves the separate fixed-unitary moment assertion; it does not
claim uniformity in a family of infinity types or replace the
finite-order principal-residue expression below.

\subsection{Keeping the full type-I mask during transfer}
For $\sigma=1/2+\delta$, define $L^T$ by removing every prime of
$T(u,d)$, not just conductor primes. The identity~\eqref{mm:masked}
is equivalently
\[
 F_{u,d}(s)=\frac{L(s,\psi_u)}{L^{T(u,d)}(2s,\psi_{u,2})}
 \prod_{p\in T(u,d),\ p\nmid\mathfrak f_u}
                         (1-\psi_u(p)(Np)^{-s}).
\]
Absolute convergence at $2\sigma>1$ bounds the reciprocal incomplete
product by $\zeta_K(2\sigma)$, without requiring any zero-free theorem
in the critical strip. For every fixed $\eta>0$,
\[
 \prod_{p\in T}(1+(Np)^{-\sigma})
 \ll_{\sigma,\eta}(N\operatorname{rad}T)^\eta,
 \qquad N\operatorname{rad}T\le C_S Nu\,Nd.
\]
Indeed $1+q^{-\sigma}\le q^\eta$ beyond a fixed threshold, and only
finitely many prime ideals lie below it. Therefore
\begin{equation}\label{mm:mask-bound}
 |F_{u,d}(\sigma+it)|
 \ll_{\nu,S,\delta,\eta}(Nu\,Nd)^\eta
                              |L(\sigma+it,\psi_u)|.
\end{equation}
This is a bounded inequality after the exact mask identity, not an
identity that removes original zeros. For example, with principal $\nu$
and $u=\pi^3$ or $\pi^6$ at $N\pi=7$, the primitive local ratio,
restoration and full masked factor at $s=1$ are respectively
\[
 \begin{array}{c|ccc}
  v_\pi(u)&\hbox{primitive ratio}&\hbox{restoration}&\hbox{masked factor}\\\hline
 3&48/49&49/48&1\\
 6&8/7&7/8&1
 \end{array}
\]
The squared primitive character in the first row is unramified at
$\pi$ although the original primitive character is ramified there.
Squaring a zero sentinel would give the wrong denominator.

For finite-order $\nu$, the normalized type-I piece is the retained
$d$ sum of the smoothed squarefree series at $Y=D/Nd$, divided by
$\sqrt D$. Its possible pole is only at $s=1$, on principal rows.
Contour shifting with the full mask gives $I_u=P_u+E_u$, with $E_u$
exactly~\eqref{mm:remainder} and
\begin{align*}
 P_u&=\mathbf1_{\psi_u=1}\sqrt D\,\widehat W(1)
 \sum_{d\ \mathrm{squarefree}\atop(d,S)=1}
 \frac{\beta_U(d)\nu(d)\chi_d(u)}{Nd}\rho_{T(u,d)},\\
 \rho_T&=\frac{\operatorname*{res}_{s=1}\zeta_K(s)}{\zeta_K(2)}
                   \prod_{p\in T}(1+(Np)^{-1})^{-1}.
\end{align*}
This does not impose positivity on $P_u$. Since $\beta_U$ is supported
on $Nd\le U^2\le D$, every inner scale $Y$ is at least one.
Four integrations by parts give
$|\widehat W(\sigma+it)|\ll_\sigma\mathcal M_4(W)(1+|t|)^{-4}$,
which justifies the shifted integrals and their horizontal limits with
the classical Hecke growth estimates.

Apply Minkowski to~\eqref{mm:remainder} in $\ell^2$ of the actual
$u$ rows, then use~\eqref{mm:mask-bound} and the off-central moment.
For sufficiently small internal losses this bounds $\|E\|_2$ by
\[
 C D^\delta H^{1/2+O(\eta)}
 \sum_d\frac{|\beta_U(d)|}{(Nd)^{\sigma-\eta}}
 \int_{\mathbb R}|\widehat W(\sigma+it)|
                     (1+|t|)^{(1+O(\eta))/2}\,dt.
\]
All original $\nu(d)\chi_d(u)$ have modulus at most one and remain
inside the exact expression before this inequality. The actual truncated
factorization gives, without replacing $\beta_U$ by arbitrary coefficients,
\[
 \sum_d\frac{|\beta_U(d)|}{(Nd)^{\sigma-\eta}}
 \le\left(\sum_{Na\le U}(Na)^{-\sigma+\eta}\right)^2
 \ll_{\delta,\eta}U^{1-2\delta+2\eta}.
\]
Choose $\eta<\delta$ if necessary and use the fixed-field ideal count.
The height integral is bounded by $C\mathcal M_4(W)$. Squaring and
absorbing the small losses in $(HD)^\varepsilon$, using $U^2\le D$,
proves~\eqref{mm:energy}. Thus the formerly conditional numerator
moment input in this residual transfer is discharged, while its
positive-power cutoff loss is retained explicitly.

\subsection{What remains open and what has been verified}
The new remainder envelope improves the old scalar-convexity $H$ power
$3/2-\delta+\varepsilon$ to $1+\varepsilon$. At fixed $U$ one may
choose $\delta$ for a requested small-loss exponent; for $U=D^a$ with
fixed $a>0$, the remaining excess tends to $2a$ as $\delta\downarrow0$.
Neither statement controls $P_u$, the other signed piece, nor their
covariance. In the original Vaughan identity the short $2a_U$ term
must be restored at smaller scales where its support can meet the
annulus. Uniform $\mathcal M_4$ bounds do not justify interchanging a
profile supremum with the row moment or assuming $U^2\le D'$ after a
scale $D'$ is reduced.

The independent prose audit is recorded in
\path{planning/analytic-oracle/arithmetic-review/adversarial-afm-review-20261104.md}.
The source-hashed finite checker tests valuation/support decompositions,
squarefull coordinates, height exponents and rational Poisson envelopes.
Its finite success does not prove reciprocity, analytic continuation,
the large sieve, the approximate functional equation or the moments.
No new Lean declarations accompany these analytic theorems. Existing
Lean local Euler identities are a separately limited verification layer.
The independent positive arithmetic determinant comparison, original
signed low-gcd short-row estimate, all required contour uniformities and
GRH remain open; Dirichlet, Hecke, Dedekind and general Artin scopes are
not identified with one another.

\paragraph{Pinned primary sources and external inputs.}
The principal analytic family input is
\href{https://arxiv.org/pdf/1112.1650v1}{Blomer--Goldmakher--Louvel,
arXiv:1112.1650v1}, Theorem~1.3, Lemma~2.1, Appendix~II and Section~5.
Our fixed-$\nu$ proof changes column coefficients and redoes the
functional-equation reduction; it does not claim the published
corollary directly states every fixed twist.
For comparison,
\href{https://arxiv.org/pdf/2201.01885v2}{Gao--Zhao,
arXiv:2201.01885v2}, Lemma~2.9 and Section~3.2, uses the sextic
large-sieve/moment interface. Its GRH-conditional central-density theorem
is not invoked. Abelian Kummer identification uses the classical
reciprocity theorem, with the project's pinned Milne \emph{Class Field
Theory}, version~4.03 (6 August 2020), Theorem~0.8 and Chapter~I,
Section~1, a reference:
\href{https://www.jmilne.org/math/CourseNotes/CFT.pdf}{Milne's notes}.
The primitive degree-one Hecke continuation and completed functional
equation are explicit classical external inputs; the project's pinned
\href{https://u.cs.biu.ac.il/~reznikov/courses/kudla-1.pdf}{Kudla notes
on Tate's thesis}, Theorem~4.3 and Corollary~4.4, supply a reference.
Their sketch and incomplete displayed epsilon formulas are not used as
independently certified constants. Standard Stirling, Mellin inversion,
Poisson representation and scalar convexity complete the external
analytic dependencies. None is replaced by a finite diagnostic or a
GRH-equivalent assumption.
